\documentclass{beamer}
\usepackage{risk231}   % shared preamble: packages, theorems, colours, listings

\setchapter{2}         % sections numbered 2.1, 2.2, ...

\title{RISK 231: Excel}

% Function signature in the style of the original slides:
% \fsig{NAME}{arguments} -> large red NAME(arguments) with italic arguments
\newcommand{\fsig}[2]{{\color{red!75!black}\large\textsf{#1(\textit{#2})}}}

\begin{document}
\maketitle

%======================================================================
\section{Intro to Excel and Functions Part 1}
%======================================================================

\sectionframe

\begin{frame}{Intro to Excel $\sim$ Tabs}
\centering
\includegraphics[width=0.85\textwidth]{ch02/ribbon_tabs}
\vspace{0.5em}
\begin{columns}[T]
\begin{column}{0.48\textwidth}
\small
\begin{itemize}
  \item \textbf{File} -- \textit{saving, printing, options}
  \item \textbf{Home} -- \textit{formatting and editing commands}
  \item \textbf{Insert} -- \textit{tables, charts, \& graphs}
  \item \textbf{Page Layout} -- \textit{headers, footers, print formats}
  \item \textbf{Formulas} -- \textit{directory of formulas}
\end{itemize}
\end{column}
\begin{column}{0.48\textwidth}
\small
\begin{itemize}
  \item \textbf{Data} -- \textit{importing data, filtering, sorting, limitations}
  \item \textbf{Review} -- \textit{proofing, commenting, protections}
  \item \textbf{View} -- \textit{spreadsheet or multi-file viewing options}
  \item \textbf{Developer} -- \textit{VBA and macros portal}
\end{itemize}
\end{column}
\end{columns}
\end{frame}

\begin{frame}{Intro to Excel $\sim$ Tabs}
Enable the ``Developer'' tab in Excel:
\begin{enumerate}
  \item \textbf{Windows version:}
  \begin{itemize}
    \item Click the ``File'' button in the top left corner of Excel.
    \item In the bottom left corner, select ``Options''.
    \item From the menu that shows up, select ``Customize Ribbon''.
    \item Check the box of the ``Developer'' tab in the right column.
  \end{itemize}
  \vspace{0.5em}
  \item \textbf{Mac version:}
  \begin{itemize}
    \item Click the ``Excel'' tab at the top of the screen.
    \item Choose the ``Preferences'' option.
    \item Select the ``Ribbon \& Toolbar'' option.
    \item Check the box to the left of ``Developer'' in the right column.
    \item Click the ``Save'' button at the bottom-right corner of the window.
  \end{itemize}
\end{enumerate}
\end{frame}

\begin{frame}{Intro to Excel $\sim$ Options setup}
\centering
\includegraphics[height=0.8\textheight]{ch02/options_formulas}
\end{frame}

\begin{frame}{Intro to Excel $\sim$ Keyboard Shortcuts}
{\Large \red{\textbf{Practice these! Know these! Use these!}}}
\vspace{1em}

A full list of keyboard shortcuts can be found here:
\begin{center}
\scriptsize
\href{https://support.office.com/en-us/article/Excel-shortcut-and-function-keys-1798d9d5-842a-42b8-9c99-9b7213f0040f}{\texttt{https://support.office.com/en-us/article/}}\\
\href{https://support.office.com/en-us/article/Excel-shortcut-and-function-keys-1798d9d5-842a-42b8-9c99-9b7213f0040f}{\texttt{Excel-shortcut-and-function-keys-1798d9d5-842a-42b8-9c99-9b7213f0040f}}
\end{center}
\end{frame}

\begin{frame}{Functions Part 1 $\sim$ Hiding/unhiding}
\small
\textbf{Hiding columns:}
\begin{itemize}
  \item Highlight the column(s) you want to hide.
  \item Right-click and select ``Hide''.
\end{itemize}
\textbf{Unhiding columns:}
\begin{itemize}
  \item Locate the hidden column(s) by finding the lightly coloured line where the
        hidden column(s) would be on the column index bar.
  \item Highlight the columns from before to after the lightly coloured line.
  \item Right-click and select ``Unhide''.
\end{itemize}
\textbf{Hiding/unhiding rows:}
\begin{itemize}
  \item Apply similar instructions for rows as given above for columns.
\end{itemize}
\end{frame}

\begin{frame}{Functions Part 1 $\sim$ Hiding/unhiding (cont.)}
\textbf{Hiding worksheets:}
\begin{itemize}
  \item Click on the worksheet tab that you would like to hide.
  \begin{itemize}
    \item Note: to hide more than one worksheet at a time, hold down Ctrl as you
          select the tabs with your mouse.
  \end{itemize}
  \item Right-click and select ``Hide''.
\end{itemize}
\textbf{Unhiding worksheets:}
\begin{itemize}
  \item Right-click on any worksheet tab.
  \item When the Unhide Worksheet window pops up, select the desired worksheet to
        unhide and click OK.
\end{itemize}
\end{frame}

\begin{frame}{Functions Part 1 $\sim$ Simple formulas \& formatting}
\textbf{Creating simple formulas:}
\begin{itemize}
  \item In the desired cell, simply write \code{=} followed by the formula. For example,
        to add the values of cells D2 and D3 into D4, simply write the formula
        \code{=D2+D3} in D4.
\end{itemize}
\vspace{0.3em}
\textbf{Formatting} (Home tab):
\begin{center}
\includegraphics[width=0.9\textwidth]{ch02/ribbon_formatting}
\end{center}
\begin{columns}[T]
\begin{column}{0.5\textwidth}
\scriptsize
\begin{itemize}
  \item \textbf{Font}: text formatting, e.g.\ bold, italics, underlined, font type and size
  \item Cell(s) border style
  \item Colour of cell(s) and font colours
\end{itemize}
\end{column}
\begin{column}{0.5\textwidth}
\scriptsize
\begin{itemize}
  \item \textbf{Alignment}: alignment of text within cells; \textit{Merge \& Center} merges
        several cells into one and centres the text; \textit{Wrap Text} makes multiple
        lines of text stay visible
  \item \textbf{Number}: number formats, e.g.\ dollar, percentage, dates, number of
        decimal places displayed, etc.
\end{itemize}
\end{column}
\end{columns}
\end{frame}

\begin{frame}{Functions Part 1 $\sim$ Format Painter}
\begin{center}
\includegraphics[height=0.9cm]{ch02/format_painter_button}
\end{center}
\begin{itemize}
  \item Clicking \textbf{once} will copy the format of cells into others.
  \vspace{0.5em}
  \item Clicking \textbf{twice} will allow you to copy the format of cells into many
        others many times until you push Esc.
\end{itemize}
\end{frame}

\begin{frame}{Functions Part 1 $\sim$ Freeze Panes}
Super useful when looking at numerous rows or columns and you still want to see the
index and/or headers.
\begin{columns}[T]
\begin{column}{0.5\textwidth}
\footnotesize
\begin{itemize}
  \item Select the cell that represents the top left of the section you want to be
        ``movable''.
  \item Go to the View tab in the ribbon and select \textbf{Freeze Panes}.
  \item Select the type of freezing you want. For example, if you want the column
        headers to stay but the rows to be fluid, select \textit{Freeze Top Row}; or, if
        you have selected cell A2, you can just select \textit{Freeze Panes}.
  \item To unfreeze panes, go to the View tab and select \textit{Unfreeze Panes}.
\end{itemize}
\end{column}
\begin{column}{0.48\textwidth}
\includegraphics[width=\textwidth]{ch02/freeze_panes}
\end{column}
\end{columns}
\end{frame}

\begin{frame}{Functions Part 1 $\sim$ ``Copy \& Paste'' vs ``Fill''}
\centering
\begin{tikzpicture}[
    card/.style={draw=black!50, rounded corners=6pt, fill=black!5,
                 minimum width=1.5cm, minimum height=1.5cm},
    st/.style={star, star points=5, star point ratio=2.3, minimum size=1.0cm,
               inner sep=0pt, draw=black!70},
    act/.style={single arrow, draw=black!60, fill=white, minimum height=2.2cm,
                single arrow head extend=0.15cm, font=\small, inner sep=3pt}]
  % copy & paste: original stays
  \node[card] (c1) at (0,0) {};
  \node[st, fill=red!80!black] at (c1) {};
  \node[act] at (2.4,0) {Copy \quad Paste};
  \node[card] (c2) at (4.8,0) {};
  \node[st, fill=red!80!black] at (c2) {};
  % cut & paste: original disappears
  \node[card, dashed, fill=none] (c3) at (0,-2) {};
  \node[st, fill=red!35, draw=black!40] at (c3) {};
  \node[act] at (2.4,-2) {Cut \quad Paste};
  \node[card] (c4) at (4.8,-2) {};
  \node[st, fill=red!80!black] at (c4) {};
\end{tikzpicture}

\vspace{0.8em}
\begin{columns}[c]
\begin{column}{0.68\textwidth}
\small
\textbf{Fill}: update numerous cells at once using either Ctrl+D or Ctrl+R, or
double-click or drag the little square (the \emph{fill handle}) at the bottom right
corner of the cell.
\end{column}
\begin{column}{0.25\textwidth}
\centering
\includegraphics[width=0.9\textwidth]{ch02/fill_handle}
\end{column}
\end{columns}
\end{frame}

\begin{frame}{Functions Part 1 $\sim$ Cell reference -- LOCAL vs ABSOLUTE}
\centering
{\Large\textbf{C3}}\\[-1pt]
{\small \red{Column \; Row}}

\vspace{0.5em}
\raggedright\small
When copying/cutting/pasting or filling, we may want the column and/or row to stay the
same for the new cell location. We can ``lock'' our references so that this can happen:
\vspace{0.3em}

\begin{tabular}{>{\color{red}\ttfamily}l p{0.78\textwidth}}
\$C\$3 & ``Absolute'' reference -- i.e.\ the reference will always be C3 \\[2pt]
C\$3   & The row is locked -- i.e.\ the reference can change columns, but the row will always refer to 3 \\[2pt]
\$C3   & The column is locked -- i.e.\ the reference can change rows, but the column will always refer to C \\[2pt]
C3     & ``Local'' reference -- i.e.\ nothing is locked and it will adjust
\end{tabular}

\vspace{0.6em}
\begin{columns}[c]
\begin{column}{0.12\textwidth}
\centering
\tikz\node[draw=black!60, rounded corners=4pt, fill=black!5, minimum size=1cm]{\Large F4};
\end{column}
\begin{column}{0.85\textwidth}
\small
When editing the formula, click on the cell reference you want to change and press
\textbf{F4} once for absolute reference, twice to lock the row, three times to lock the
column, and four times for local reference.
\end{column}
\end{columns}
\end{frame}

%======================================================================
\section{Excel: Functions Part 2}
%======================================================================

\sectionframe

\begin{frame}{Functions Part 2 $\sim$ Functions}
\small
Excel has hundreds of built-in functions that you can use in your calculations and data
wrangling. You provide the name of the function and the required (and/or optional) inputs
and Excel will run the code behind the scenes to provide your output.

\vspace{0.5em}
You will follow this structure:
\begin{center}
\red{\large NAME of FUNCTION (\textit{necessary input, [optional input]})}
\end{center}

You can also nest functions inside functions. For example, if you wanted to use a
particular function to calculate an output that would be used as an input into another
function, you could do it all in one step by nesting the functions:
\begin{center}
\red{NAME of FUNCTION (}\green{\textit{NAME of FUNCTION (necessary input, [optional input])}}\red{\textit{, [optional input]})}
\end{center}

For example:
\begin{center}
\large\code{=\red{SUM(}\green{AVERAGE(D3:G3)}\red{,} \green{AVERAGE(A3:C3)}\red{)}}
\end{center}
\end{frame}

\begin{frame}{Functions Part 2 $\sim$ VLOOKUP}
\begin{center}
\fsig{VLOOKUP}{lookup\_value, table\_array, col\_index\_num, [range\_lookup]}
\end{center}
\vspace{0.3em}
\begin{tabular}{>{\ttfamily\color{red!75!black}}l p{0.62\textwidth}}
\toprule
lookup\_value   & What are you looking up? \\
table\_array    & What is the range of cells where your data lives for the lookup? \\
col\_index\_num & Column number of the ranged data that you want to return when the
                  lookup value is found for that row \\
{[range\_lookup]} & Do you want an exact match (\code{FALSE}) or do you want Excel to
                  guess (\code{TRUE}, approximate match)? \\
\bottomrule
\end{tabular}

\vspace{1em}
\textbf{Important:} \textit{Remember that the lookup value should always be in the
\textbf{first column} of the range for VLOOKUP to work correctly.}

\vspace{0.5em}
e.g.\ \code{=VLOOKUP(A2,B:C,2,FALSE)}
\end{frame}

\begin{frame}{Functions Part 2 $\sim$ IF}
\begin{center}
\fsig{IF}{logical\_test, [value\_if\_true], [value\_if\_false]}
\end{center}
\begin{tabular}{>{\ttfamily\color{red!75!black}}l p{0.6\textwidth}}
\toprule
logical\_test      & The requirement that leads to the THEN and ELSE steps \\
{[value\_if\_true]}  & What happens if the logical test is true? \\
{[value\_if\_false]} & What happens if the logical test is false? \\
\bottomrule
\end{tabular}

\vspace{0.4em}
\small
i.e.\ \emph{If something is true, then do something, otherwise, do a different thing\ldots}

\vspace{0.6em}
\textbf{Nested IFs} are frequently used. For example, if something is true, then do
something, otherwise, if something else is true, then do something, otherwise\ldots
\begin{center}
\red{IF (\textit{logical test, [value if true],}} \green{IF (\textit{logical test, [value if true],}}
\purple{IF (\textit{logical test, [value if true], [value if false]})}\green{)}\red{)}
\end{center}
\end{frame}

\begin{frame}{Functions Part 2 $\sim$ Named Ranges}
\footnotesize
To make writing functions and referencing data sets in general easier, Excel allows you
to name selected data ranges so that you don't have to list the cell references every
time they are used.

\vspace{0.4em}
For example, if you have data in cells A1:G23, you can name the range ``my\_data'' and
simply refer to it by its name. Easy-peasy!

\vspace{0.4em}
\textit{Two ways to name a range:}
\vspace{0.3em}

\begin{columns}[c]
\begin{column}{0.58\textwidth}
\begin{itemize}
  \item Highlight the region and enter the name in the \textbf{Name Box} located in the
        top left corner of the worksheet, just below the ribbon, and then press Enter.
\end{itemize}
\end{column}
\begin{column}{0.4\textwidth}
\includegraphics[width=0.85\textwidth]{ch02/name_box}
\end{column}
\end{columns}
\vspace{0.4em}
\begin{columns}[c]
\begin{column}{0.58\textwidth}
\begin{itemize}
  \item Go to the Formulas tab of the ribbon and click on \textbf{Name Manager}. Select
        New, name the range and define the range in the pop-up. You can refer back to
        the Name Manager to update, add, or delete your named ranges.
\end{itemize}
\end{column}
\begin{column}{0.4\textwidth}
\includegraphics[width=0.65\textwidth]{ch02/name_manager}
\end{column}
\end{columns}
\end{frame}

\begin{frame}{Functions Part 2 $\sim$ Data Validation}
\small
Not sure why they call it Data Validation\ldots

If you would like to use a dropdown to make things a bit easier for yourself and any user
of your Excel file, follow the steps below. There are other really cool options with Data
Validation as well.
\begin{itemize}
  \item Make a list of your desired selections somewhere within your workbook.
  \item Select the cell for which you would like a dropdown.
  \item On the Data tab of the ribbon, select \textbf{Data Validation}.
\end{itemize}
\begin{center}
\includegraphics[width=0.7\textwidth]{ch02/data_validation_ribbon}
\end{center}
\begin{itemize}
  \item Select \textbf{List} and put the range of cells where your desired selections
        live, and then select OK.
\end{itemize}
\end{frame}

\begin{frame}{Functions Part 2 $\sim$ INDIRECT}
\begin{center}
\fsig{INDIRECT}{ref\_text, [a1]}
\end{center}
\vspace{1em}
The INDIRECT function allows you to use the \emph{contents} of a cell as a reference,
rather than the cell itself.
\begin{itemize}
  \item i.e.\ When used as a nested function, it can grab the contents of a cell and use
        it as an input into another function.
\end{itemize}
\vspace{0.5em}
This can be particularly useful when referring to worksheet tab names, file names, and
named ranges.

\vspace{0.5em}
e.g.\ if cell A1 contains the text \code{my\_data}, then \code{=SUM(INDIRECT(A1))} sums the
named range \code{my\_data}.
\end{frame}

%======================================================================
\section{Excel: Functions Part 3}
%======================================================================

\sectionframe

\begin{frame}{Functions Part 3 $\sim$ INDEX}
\begin{center}
\fsig{INDEX}{array, row\_num, [column\_num]}
\end{center}
\begin{tabular}{>{\ttfamily\color{red!75!black}}l p{0.62\textwidth}}
\toprule
array          & Range of data you are using \\
row\_num       & Relative row number within the grid of selected data that you are
                 trying to identify \\
{[column\_num]} & Relative column number within the grid of selected data that you are
                 trying to identify \\
\bottomrule
\end{tabular}

\vspace{0.8em}
\small
The INDEX function is very useful for larger datasets and in conjunction with other
functions. Usually, you'll be using other functions as inputs into the INDEX function.

\vspace{0.5em}
Most of the time, you'll use VLOOKUP instead because it is more efficient from a memory
usage perspective, but there will be times when INDEX is necessary when VLOOKUP is less
handy.
\end{frame}

\begin{frame}{Functions Part 3 $\sim$ OFFSET and MATCH}
\begin{center}
\fsig{OFFSET}{reference, rows, columns}
\end{center}
\begin{tabular}{>{\ttfamily\color{red!75!black}}l p{0.6\textwidth}}
\toprule
reference & Starting (origin) cell \\
rows      & How many rows down from the origin? \\
columns   & How many columns to the right from the origin? \\
\bottomrule
\end{tabular}

\vspace{0.6em}
\begin{center}
\fsig{MATCH}{lookup\_value, lookup\_array, [match\_type]}
\end{center}
\begin{tabular}{>{\ttfamily\color{red!75!black}}l p{0.6\textwidth}}
\toprule
lookup\_value   & What am I trying to match? \\
lookup\_array   & What range of cells am I looking into for my lookup value? \\
{[match\_type]} & Do I want an exact match \code{0} or an approximate match \code{1}? \\
\bottomrule
\end{tabular}

\vspace{0.5em}
\small
MATCH will return the reference \emph{position number}. Super useful within other
functions, like VLOOKUP, INDEX, and OFFSET, e.g.\
\code{=INDEX(B2:E20, MATCH("Bako",A2:A20,0), 3)}.
\end{frame}

\begin{frame}{Functions Part 3 $\sim$ Text to Columns}
\small
Many times, to save on space, databases will store data in simpler formats, like comma
separated values (csv). In order to take the raw data and translate it into a grid format
like we have in Excel, we can use the \textbf{Text to Columns} tool.
\vspace{0.3em}
\begin{columns}[c]
\begin{column}{0.52\textwidth}
\begin{itemize}
  \item After pasting (or inserting) your data into the worksheet, oftentimes all in
        column A, highlight the data.
  \item After highlighting, go to the Data tab on the ribbon and select
        \textit{Text to Columns}.
\end{itemize}
\end{column}
\begin{column}{0.46\textwidth}
\includegraphics[width=\textwidth]{ch02/text_to_columns_ribbon}
\end{column}
\end{columns}
\end{frame}

\begin{frame}{Functions Part 3 $\sim$ Text to Columns (cont.)}
\small
\begin{columns}[c]
\begin{column}{0.66\textwidth}
\begin{itemize}
  \item If your data is separated with special characters, like commas, choose
        ``Delimited''. If you want to manually select how the data is split, choose
        ``Fixed width''. Click Next.
\end{itemize}
\end{column}
\begin{column}{0.3\textwidth}
\includegraphics[width=\textwidth]{ch02/text_to_columns_delimited}
\end{column}
\end{columns}
\vspace{0.3em}
\begin{columns}[c]
\begin{column}{0.66\textwidth}
\begin{itemize}
  \item Select the special character that separates your data. You'll see a preview of
        how the data will be split below. Click Next. On the next screen, you will
        almost always select ``General'' and then Finish.
\end{itemize}
\end{column}
\begin{column}{0.3\textwidth}
\centering
\includegraphics[width=0.75\textwidth]{ch02/text_to_columns_delimiters}
\end{column}
\end{columns}
\vspace{0.3em}
\begin{itemize}
  \item If you want to manually select how the data is split, after choosing ``Fixed
        width'', the next screen will allow you to make your column split selections.
\end{itemize}
\end{frame}

\begin{frame}{Functions Part 3 $\sim$ Goal Seek}
\small
Sometimes, you'll have several pieces of data, but you need to find another piece by
converging on an answer by allowing some other variable to fluctuate. This situation
comes up quite frequently in Finance and as an Actuary. We can use the \textbf{Goal Seek}
tool to do this.
\begin{columns}[c]
\begin{column}{0.6\textwidth}
\begin{itemize}
  \item Go to the Data tab on the ribbon and select \textit{What-If Analysis}. In the
        dropdown, select \textit{Goal Seek}.
\end{itemize}
\end{column}
\begin{column}{0.38\textwidth}
\includegraphics[width=0.8\textwidth]{ch02/goal_seek_menu}
\end{column}
\end{columns}
\vspace{0.2em}
\begin{columns}[c]
\begin{column}{0.4\textwidth}
\includegraphics[width=0.9\textwidth]{ch02/goal_seek_dialog}
\end{column}
\begin{column}{0.58\textwidth}
\footnotesize
\begin{description}
  \item[Set cell:] Which cell am I fixing to a certain value?
  \item[To value:] What value am I aiming for? Note that this is usually zero because we
        are often using Goal Seek for convergence purposes. It can be anything you want,
        though.
  \item[By changing cell:] Which cell is allowed to fluctuate?
\end{description}
\end{column}
\end{columns}
\end{frame}

%======================================================================
\section{Excel: Loan Calculator}
%======================================================================

\sectionframe

\begin{frame}{Amortization}
\begin{itemize}
  \item When you have a large expense (such as a home or car purchase), it would be quite
        rare that you purchase it only with cash, so we rely on available loans to
        enable such purchases.
  \vspace{0.5em}
  \item Once a loan is initiated, it should be repaid within the pre-determined term by
        making periodic payments (weekly, bi-weekly, or monthly).
  \vspace{0.5em}
  \item Such a process is called `A\red{mort}ization' -- \red{killing off} the loan over
        time (recall the Latin root word `\red{mort}' -- which means death; it also
        appears in \red{mort}gage).
\end{itemize}
\end{frame}

\begin{frame}{Amortization Schedule}
\begin{itemize}
  \item There are four main factors that interact and determine the amortization
        schedule:
  \begin{itemize}
    \item Annual percentage rate (APR)
    \item Length of the term
    \item Loan principal
    \item Periodic payment amount
  \end{itemize}
  \vspace{0.5em}
  \item If any three of the above are given, then the remaining item is also determined
        by the amortization schedule (this requires Goal Seek to solve for it).
\end{itemize}
\end{frame}

\begin{frame}{Interactions among the factors}
\begin{itemize}
  \item All else being equal, the periodic payment amount would increase if
  \begin{itemize}
    \item the APR increases,
    \item the loan principal increases, and
    \item the length of the term decreases.
  \end{itemize}
  \vspace{0.5em}
  \item Other than Excel, there are many available online calculators that offer an
        amortization schedule:
        \url{https://www.calculator.net/finance-calculator.html}
  \vspace{0.5em}
  \item More details are available in ACMA 201 (time value of money, mortgages, bonds,
        etc.)
\end{itemize}
\end{frame}

\begin{frame}{Example: Himchan's mortgage}
\centering
\includegraphics[width=0.95\textwidth]{ch02/mortgage_example}

\vspace{0.6em}
A mortgage needs to be renewed every 3 or 5 years in Canada!\\
(mine is 5 years fixed-rate -- see the \textbf{Maturity Date} and
\textbf{Maturity Balance})
\end{frame}

%======================================================================
\section{Excel: Functions Part 4}
%======================================================================

\sectionframe

\begin{frame}{Functions Part 4 $\sim$ IFERROR}
\begin{center}
\fsig{IFERROR}{value, value\_if\_error}
\end{center}
\begin{tabular}{>{\ttfamily\color{red!75!black}}l p{0.62\textwidth}}
\toprule
value           & Doing something, usually a function (e.g.\ VLOOKUP) \\
value\_if\_error & What you want the output to be if there is an error with the function
                  (e.g.\ set equal to zero) \\
\bottomrule
\end{tabular}

\vspace{0.8em}
i.e.\ \emph{``If this action results in an error, then use this as the output, otherwise,
use the output from the action''}

\vspace{0.8em}
e.g.\ \code{=IFERROR(VLOOKUP(A2,B:C,2,FALSE),0)}
\end{frame}

\begin{frame}{Functions Part 4 $\sim$ SUMIF}
\begin{center}
\fsig{SUMIF}{range, criteria, sum\_range}
\end{center}
\begin{tabular}{>{\ttfamily\color{red!75!black}}l p{0.66\textwidth}}
\toprule
range      & Define the range to which the criteria applies \\
criteria   & Define the criteria that triggers items to be included in the sum \\
sum\_range & Define the range of data to be summed if the criteria is met in the given
             range for the first input \\
\bottomrule
\end{tabular}

\vspace{0.8em}
i.e.\ \emph{``If the criteria is true for cells in a given range, then sum the values in
the same rows in another (or the same) given range''}

\vspace{0.8em}
e.g.\ \code{=SUMIF(A2:A100,"CA",C2:C100)}
\end{frame}

\begin{frame}{Functions Part 4 $\sim$ Filtering and Sorting}
\small
\begin{columns}[c]
\begin{column}{0.6\textwidth}
{\large\red{Filtering}}\\\footnotesize
Select any cell in the column header of your desired table, then use your keyboard
shortcut (Ctrl+Shift+L) or click on ``Filter'' in the Data ribbon tab. To use the filter,
open the dropdowns.
\end{column}
\begin{column}{0.38\textwidth}
\includegraphics[width=\textwidth]{ch02/filter_button}
\end{column}
\end{columns}
\vspace{0.2em}
{\large\red{Sorting}}\\\footnotesize
There are several ways to get to sorting capabilities: using the dropdown of the filter in
your table, using the ``Sort'' buttons on the Data ribbon tab, or highlighting a range of
cells, right-clicking and selecting ``Sort''.

\vspace{0.3em}
Usually, you will want to sort where the rest of a table adjusts accordingly (keeping the
rows the same, but in a different order). Sometimes, however, you may want to only adjust
a portion of the data, leaving everything else in the row the same. You would use
\textbf{custom sort} for this.
\begin{center}
\includegraphics[width=0.62\textwidth]{ch02/sort_dialog}
\end{center}
\end{frame}

\begin{frame}{Functions Part 4 $\sim$ VALUE, TEXT, LEFT, RIGHT, MID}
\small
\fsig{VALUE}{text}\\
Turns the given text-filled cell into a value, if possible.

\vspace{0.4em}
\fsig{TEXT}{value, format\_text}\\
Turns the given value-filled cell into text, given the specific text format desired. There
are A LOT of different text formats (e.g.\ \code{"0,000.00"} or \code{"mm/dd/yyyy"}); see
Excel's help page for TEXT for further examples.

\vspace{0.4em}
\fsig{LEFT}{text, num\_chars}\\
Returns the given number of characters from a cell, counting from the left (be careful of
leading spaces!)

\vspace{0.4em}
\fsig{RIGHT}{text, num\_chars}\\
Returns the given number of characters from a cell, counting from the right (be careful
of trailing spaces!)

\vspace{0.4em}
\fsig{MID}{text, start\_num, num\_chars}\\
Returns the given number of characters from a cell, counting from a given starting
position. (e.g.\ if A5 contains 1234567, then \code{MID(A5,3,2)} = 34)
\end{frame}

\begin{frame}{Functions Part 4 $\sim$ DATE, YEAR, MONTH, and DAY}
\fsig{DATE}{year, month, day}\\
Creates a value in date format, given year, month, and day values. These could be
manually given or pulled from other portions of cell values using other functions.

\vspace{0.6em}
\fsig{YEAR}{serial\_number}\\
Returns the year from a date format, which is stored in the Excel metadata as a serial
number (i.e.\ the date format displays this serial number in a format we are used to
seeing).

\vspace{0.6em}
\fsig{MONTH}{serial\_number}\\
Returns the month number from a date format.

\vspace{0.6em}
\fsig{DAY}{serial\_number}\\
Returns the day number from a date format.
\end{frame}

\begin{frame}{Functions Part 4 $\sim$ FIND and LEN}
\begin{center}
\fsig{FIND}{find\_text, within\_text, [start\_num]}
\end{center}
\begin{tabular}{>{\ttfamily\color{red!75!black}}l p{0.64\textwidth}}
\toprule
find\_text    & What you are looking for in the string -- the FIND function will return
                its position number \\
within\_text  & Location of the string \\
{[start\_num]} & Starting position in the string from which to search (most of the time you
                will not enter this because you most often start at the beginning of the
                string) \\
\bottomrule
\end{tabular}

\vspace{0.5em}
i.e.\ \emph{``The FIND function will return the position number of a particular text you
are looking for in a given text string''}

\vspace{1em}
\fsig{LEN}{text} \quad Returns the numerical length of a given string.
\end{frame}

%======================================================================
\section{Excel: Pivot Tables}
%======================================================================

\sectionframe

\begin{frame}{Pivot Tables}
Pivot tables are a fast, flexible, and user-friendly way to summarize and analyze data in
Excel. They are particularly efficient for ad-hoc analysis.
\begin{itemize}
  \item Reference data should be organized as a list or table
  \item All columns should have headers
  \item Can't have any empty rows or columns
\end{itemize}
\begin{center}
\footnotesize
\begin{tabular}{lrr}
\toprule
\textbf{Row Labels} & \textbf{Count of PolNum} & \textbf{Sum of Face Amt} \\
\midrule
Alfonzo  & 26 & 5236448 \\
Bako     & 40 & 5492664 \\
Barajas  & 44 & 6289588 \\
Barrett  & 53 & 7846295 \\
Bennett  & 37 & 5132155 \\
Betts    & 40 & 4893288 \\
Blanco   & 46 & 5784722 \\
Bowen    & 40 & 6246566 \\
\midrule
\textbf{Grand Total} & \textbf{326} & \textbf{46921726} \\
\bottomrule
\end{tabular}
\end{center}
\end{frame}

\begin{frame}{Pivot Tables $\sim$ Getting started}
\small
\begin{columns}[c]
\begin{column}{0.58\textwidth}
To get started, select the starting cell where you want the pivot table to be placed
(i.e.\ where the top left corner of the pivot will be placed on the sheet). Then, click
the \textbf{PivotTable} button on the ``Insert'' ribbon tab.
\end{column}
\begin{column}{0.38\textwidth}
\includegraphics[width=\textwidth]{ch02/pivot_insert_button}
\end{column}
\end{columns}
\vspace{1em}
\begin{columns}[c]
\begin{column}{0.5\textwidth}
\includegraphics[width=\textwidth]{ch02/pivot_create_dialog}
\end{column}
\begin{column}{0.48\textwidth}
Click on the arrow button (at the right of \textit{Table/Range}) and select the data you
would like your pivot table to be based on, making sure that the column headers of the
data are the first row in your selected data range. Click the arrow again and select OK.

\textit{\blue{(e.g.\ Sales Data 2)}}
\end{column}
\end{columns}
\end{frame}

\begin{frame}{Pivot Tables $\sim$ The Field List}
\begin{columns}[c]
\begin{column}{0.7\textwidth}
\scriptsize
The \textbf{PivotTable Field List} is where you build the structure of the pivot.
\begin{itemize}
  \item \textbf{Field list:} the column headers from your data source will be listed here.
  \item \textbf{\blue{Rows:}} the desired rows for your pivot. Just drag a column header
        from the list and place it in this box. You can have more than one category here
        as well, but this may complicate your pivot table. \textit{(e.g.\ Gender and Class)}
  \item \textbf{\red{Values:}} what data do you want to aggregate by the row splits?
        Examples could be things like aggregated premiums or claim counts.
        \textit{(e.g.\ Face Amt)}
  \item \textbf{\textcolor{cyan!60!black}{Columns:}} you may want to split the data into
        even finer groupings. You can add additional columns by dragging column headers
        from the list and placing them in this box. You can have more than one here, but
        this may complicate your pivot table. \textit{(e.g.\ PremMode)}
  \item \textbf{\orange{Filters:}} you can also add one or many filters by dragging column
        headers here. \textit{(e.g.\ drag Gender from the Rows box to the Filters box)}
\end{itemize}
\end{column}
\begin{column}{0.28\textwidth}
\centering
\begin{tikzpicture}
  \node[inner sep=0, anchor=south west] (img) at (0,0)
       {\includegraphics[height=0.82\textheight]{ch02/pivot_field_list}};
  \begin{scope}[x={(img.south east)}, y={(img.north west)}, very thick]
    \draw[black!70]       (0.02,0.585) rectangle (0.98,0.875);
    \draw[orange]         (0.025,0.365) rectangle (0.495,0.555);
    \draw[cyan!60!black]  (0.505,0.365) rectangle (0.975,0.555);
    \draw[blue]           (0.025,0.035) rectangle (0.495,0.355);
    \draw[red]            (0.505,0.035) rectangle (0.975,0.355);
  \end{scope}
\end{tikzpicture}
\end{column}
\end{columns}
\end{frame}

\begin{frame}{Pivot Tables $\sim$ Value Field Settings}
\small
\begin{columns}[T]
\begin{column}{0.48\textwidth}
We can choose many different aggregation methods (Sum, Count, Average, Max, Min,
Product, \ldots) for the data in the body of our Pivot Table, by the different column
content.

\vspace{0.4em}
Either right-click on the Pivot Table column or click on the field in the Field List,
and select ``Value Field Settings''. You can see that we have MANY choices here.

\vspace{0.4em}
If you click on ``Number Format'', you can also select the way the data is displayed in
the Pivot Table.
\end{column}
\begin{column}{0.5\textwidth}
\centering
\includegraphics[height=0.33\textheight]{ch02/pivot_value_field_menu}\hfill
\includegraphics[height=0.33\textheight]{ch02/pivot_value_field_settings}\\[4pt]
\includegraphics[height=0.33\textheight]{ch02/pivot_number_format}
\end{column}
\end{columns}
\end{frame}

\begin{frame}{Pivot Tables $\sim$ PivotTable Tools: Analyze}
\small
When clicking anywhere in your Pivot Table, not only do we see the Field List, but the
ribbon adds a new tab for \textbf{PivotTable Tools}.
\begin{center}
\includegraphics[height=1.3cm]{ch02/pivot_tools_analyze}
\end{center}
The \textbf{Analyze} tab contains many capabilities.
\begin{itemize}
  \item All the way to the left, you can name your Pivot Table.
  \item You can also access the PivotTable Options, which have MANY different choices of
        settings for your table.
  \item The Analyze tab also has the ability to make \textbf{PivotCharts}, which are built
        directly off of your Pivot Table and will refresh when you refresh your table.
  \item Within ``Fields, Items, \& Sets'', you can create \textbf{calculated fields},
        which are often helpful.
  \item You can make the Field List disappear and reappear by clicking on the button
        located on the right of the Analyze tab.
\end{itemize}
\end{frame}

\begin{frame}{Pivot Tables $\sim$ Refresh}
\small
When you set the data source for your Pivot Table, a copy of that data source is created
in the metadata (i.e.\ ``behind the scenes''). The Pivot Table is actually pulling from
this \textbf{Pivot Table Cache}. This allows you to make changes in the data without
directly affecting your Pivot Tables until you desire.

\vspace{0.4em}
As such, if you make changes to your data and you want those changes reflected in the
Pivot Table, you will need to click on the \textbf{Refresh} button, which updates the
Pivot Table Cache with the new data.
\vspace{0.6em}
\begin{center}
\includegraphics[height=2.4cm]{ch02/pivot_refresh_data_tab}
\quad\raisebox{1.1cm}{\red{\textit{or}}}\quad
\includegraphics[height=2.4cm]{ch02/pivot_refresh_analyze_tab}
\end{center}
\end{frame}

\begin{frame}{Pivot Tables $\sim$ Change Data Source}
You may also want to change the data used for your Pivot Table (e.g.\ maybe you like the
structure of the Pivot and have another data source with the same column headers, but
newer data).

\vspace{0.5em}
Use \textbf{Change Data Source} on the Analyze tab:
\begin{center}
\includegraphics[height=3cm]{ch02/pivot_refresh_analyze_tab}
\end{center}
\vspace{0.3em}
\textit{\blue{(e.g.\ Change the data to Sales Data 1)}}
\end{frame}

\begin{frame}{Pivot Tables $\sim$ PivotTable Tools: Design}
When clicking anywhere in your Pivot Table, not only do we see the Field List, but the
ribbon adds a new tab for \textbf{PivotTable Tools}.
\begin{center}
\includegraphics[height=1.3cm]{ch02/pivot_tools_design}
\end{center}
On the \textbf{Design} tab, you can change the colour and format of your Pivot Table. You
can also decide whether or not to have row and/or column headers, and how/whether to
display subtotals and/or grand totals.

\vspace{1em}
\textit{\blue{(e.g.\ Change our Pivot Table to another colour without Grand Totals)}}
\end{frame}

\begin{frame}{Pivot Tables $\sim$ Drilling down}
If you would like to view the specific underlying data for a particular value seen in
your Pivot Table, simply \textbf{double-click} on the value and a new worksheet will be
created with the rows associated with this value.

\vspace{0.5em}
For example, maybe you noticed in your Pivot Table that for Males in the state of
California, the premiums look a little strange for Territory 5. You can double-click on
this value in the Pivot Table and it will supply all the rows of data for California
Males in Territory 5 from your data source. Super helpful for deep dive analysis!

\vspace{1em}
\textit{\blue{(e.g.\ Try it out on our table by double-clicking on any chosen value)}}
\end{frame}

%======================================================================
\section{Excel: Charts and Conditional Formatting}
%======================================================================

\sectionframe

\begin{frame}{Charts}
Charts (graphs) are another way to summarize data in Excel. The advantage of using charts
is that they allow you to present information graphically, which can help you or your
audience visualize trends or comparisons.
\vspace{0.5em}
\begin{columns}[T]
\begin{column}{0.48\textwidth}
There are many different types of charts built into Excel. We will focus on the following
most common types:
\begin{itemize}
  \item Bar Charts
  \item Line Charts
  \item Pie Charts
\end{itemize}
\end{column}
\begin{column}{0.48\textwidth}
Components of a Chart:
\begin{itemize}
  \item Vertical and Horizontal Axes
  \item Data Series
  \item Legend
  \item Chart Title
  \item Axis Titles
\end{itemize}
\end{column}
\end{columns}
\end{frame}

\begin{frame}{Charts $\sim$ Building a chart}
\small
\begin{columns}[c]
\begin{column}{0.45\textwidth}
To build charts, go to the \textbf{Insert} tab on the ribbon and look for the Charts
section in the middle of the ribbon. The first row of icons holds the \textbf{bar
charts}, the second row the \textbf{line charts}, and below them the \textbf{pie
charts}.
\end{column}
\begin{column}{0.5\textwidth}
\includegraphics[width=\textwidth]{ch02/chart_ribbon}
\end{column}
\end{columns}
\vspace{0.6em}
\begin{columns}[c]
\begin{column}{0.4\textwidth}
Select the type of chart you would like, then select ``Select Data'' on the ribbon tab
that shows.

\vspace{0.4em}
Select the data desired. It should populate the axes properly, but if not, you can change
your selections here or later.
\end{column}
\begin{column}{0.58\textwidth}
\includegraphics[width=\textwidth]{ch02/chart_select_data}
\end{column}
\end{columns}
\end{frame}

\begin{frame}{Charts $\sim$ Chart buttons}
\small
By clicking in the chart, you will see the following buttons show to the right of the
chart:
\vspace{0.3em}
\begin{columns}[c]
\begin{column}{0.1\textwidth}
\includegraphics[height=3cm]{ch02/chart_side_buttons}
\end{column}
\begin{column}{0.86\textwidth}
\begin{itemize}
  \item[\textbf{+}] Click on this to add items to the chart. For example, you may want
        to add a legend, axis titles, or other items.
  \vspace{0.5em}
  \item[\textbf{Brush}] Click on this to change the formatting of your chart. There are
        many options here to make your chart look any way you'd like.
  \vspace{0.5em}
  \item[\textbf{Filter}] Click on this to filter the data you'd like to see in your chart.
\end{itemize}
\end{column}
\end{columns}
\vspace{0.5em}
\begin{columns}[c]
\begin{column}{0.4\textwidth}
Also note that you can easily change the data ranges underlying your chart by adjusting
the coloured sections showing on your data when clicking in the chart.
\end{column}
\begin{column}{0.56\textwidth}
\includegraphics[width=\textwidth]{ch02/chart_data_ranges}
\end{column}
\end{columns}
\end{frame}

\begin{frame}{Charts $\sim$ Chart Tools}
By clicking in the chart, you will also see a \textbf{Chart Tools} tab added to the
ribbon with the following sub-tabs:
\begin{center}
\includegraphics[height=1.4cm]{ch02/chart_tools_tabs}
\end{center}
\begin{itemize}
  \item The \textbf{Design} tab will allow you to change the chart type, change the way
        the chart looks, change the data, among other things.
  \item The \textbf{Format} tab allows you to change the chart titles, fonts, colours,
        etc.
  \item You can also get to formatting by double-clicking on the chart -- a side window
        will appear with many formatting options.
\end{itemize}
\end{frame}

\begin{frame}{Conditional Formatting}
\small
\begin{columns}[c]
\begin{column}{0.6\textwidth}
In the Home tab of the ribbon, there is a choice for \textbf{Conditional Formatting},
which will give you the option of formatting cells based on their contents and other
possible rules.

\vspace{0.6em}
Just highlight the cells you want to be conditionally formatted and then select the type
of formatting and criteria desired.

\vspace{0.6em}
You can also manage existing rules for previously assigned formatting. A \emph{rule} is
the criteria given and the resultant formatting.

\vspace{0.6em}
There are other cool things you can do, like small data bars or heat mapping the cells.

\vspace{0.6em}
Cool stuff!
\end{column}
\begin{column}{0.36\textwidth}
\centering
\includegraphics[height=0.8\textheight]{ch02/conditional_formatting_menu}
\end{column}
\end{columns}
\end{frame}

%======================================================================
\section{Excel: Functions Part 5}
%======================================================================

\sectionframe

\begin{frame}{Functions Part 5 $\sim$ PMT, RATE, and NPER}
\fsig{PMT}{rate, nper, pv, [fv], [type]}\\
i.e.\ \emph{``Calculates the payment for a loan (or annuity) based on constant payments
and a constant interest rate''}

\vspace{0.8em}
\fsig{RATE}{nper, pmt, pv, [fv], [type], [guess]}\\
i.e.\ \emph{``Returns the interest rate per period of a loan or an investment''}
\begin{itemize}
  \item \code{guess}: giving Excel a ``starting point'' to figure out the interest rate.
        You can usually just enter a ballpark number, like 0.01.
\end{itemize}

\vspace{0.5em}
\fsig{NPER}{rate, pmt, pv, [fv], [type]}\\
i.e.\ \emph{``Returns the number of periods for an investment based on periodic, constant
payments and a constant interest rate''}
\end{frame}

\begin{frame}{Functions Part 5 $\sim$ Print Preparation}
\centering
View tab $\rightarrow$ \textbf{Page Break Preview}\\[3pt]
\includegraphics[height=2.3cm]{ch02/page_break_preview}

\vspace{1em}
Page Layout tab $\rightarrow$ \textbf{Margins}, \textbf{Orientation}, \textbf{Print Area},
and the \textbf{Page Setup} dialog launcher (bottom-right corner)\\[3pt]
\includegraphics[height=2.3cm]{ch02/page_setup_ribbon}
\end{frame}

\begin{frame}{Functions Part 5 $\sim$ Print Preparation: Page Setup}
\centering
\small Page Layout $\rightarrow$ Page Setup dialog launcher $\rightarrow$ \textbf{Page} tab (orientation, scaling, paper size)\\[4pt]
\includegraphics[height=0.74\textheight]{ch02/page_setup_page}
\end{frame}

\begin{frame}{Functions Part 5 $\sim$ Print Preparation: Page Setup (cont.)}
\centering
\small Page Layout $\rightarrow$ Page Setup dialog launcher $\rightarrow$ \textbf{Margins} tab\\[4pt]
\includegraphics[height=0.74\textheight]{ch02/page_setup_margins}
\end{frame}

\begin{frame}{Functions Part 5 $\sim$ Print Preparation: Page Setup (cont.)}
\centering
\small Page Layout $\rightarrow$ Page Setup dialog launcher $\rightarrow$ \textbf{Header/Footer} tab\\[4pt]
\includegraphics[height=0.74\textheight]{ch02/page_setup_header_footer}
\end{frame}

\begin{frame}{Functions Part 5 $\sim$ Print Preparation: Page Setup (cont.)}
\centering
\small Page Layout $\rightarrow$ Page Setup dialog launcher $\rightarrow$ \textbf{Sheet} tab (print area, rows to repeat at top, page order)\\[4pt]
\includegraphics[height=0.74\textheight]{ch02/page_setup_sheet}
\end{frame}

\begin{frame}{Functions Part 5 $\sim$ Print Preparation: Print Preview}
\begin{columns}[c]
\begin{column}{0.5\textwidth}
\red{Print Preview}
\vspace{0.6em}
\begin{itemize}
  \item File $\rightarrow$ Print
  \vspace{0.4em}
  \item[] Or
  \vspace{0.4em}
  \item Page Layout $\rightarrow$ Page Setup $\rightarrow$ Print Preview
\end{itemize}
\end{column}
\begin{column}{0.45\textwidth}
\centering
\includegraphics[height=0.82\textheight]{ch02/print_backstage}
\end{column}
\end{columns}
\end{frame}

\begin{frame}{Functions Part 5 $\sim$ Some Additional Formatting}
\footnotesize
\begin{columns}[c]
\begin{column}{0.68\textwidth}
\textbf{Wrap Text}: \textit{wrap extra long text into multiple lines within the same cell
so that you can see all of it.}

\vspace{0.4em}
\textbf{Merge \& Center}: \textit{combine and centre the contents of the selected cells in a
new larger cell. This is a great way to create a label that spans multiple columns.}
\end{column}
\begin{column}{0.28\textwidth}
\includegraphics[width=\textwidth]{ch02/wrap_merge_buttons}
\end{column}
\end{columns}
\vspace{0.4em}
\begin{columns}[c]
\begin{column}{0.68\textwidth}
\textbf{Format Painter}: \textit{lets you copy all of the formatting from one group of cells
and apply it to another group. Think of it as copying and pasting for formatting.}
\end{column}
\begin{column}{0.28\textwidth}
\includegraphics[width=\textwidth]{ch02/format_painter_clipboard}
\end{column}
\end{columns}
\vspace{0.4em}
\begin{columns}[c]
\begin{column}{0.6\textwidth}
\textbf{Column Widths}: \textit{you can change the widths of multiple columns uniformly by
highlighting the desired columns and then adjusting the width of the first column in your
selection with your mouse. This will update all selected columns to be the same width.}
\end{column}
\begin{column}{0.36\textwidth}
\includegraphics[width=\textwidth]{ch02/column_width}
\end{column}
\end{columns}
\end{frame}

\end{document}
